\chapter{Conclusions}
\label{c:conclusions}

Identifying the origin of an epidemic outbreak from a snapshot of the infection
state is a problem of both scientific and practical importance. On temporal
contact networks, the task is complicated by the fact that transmission
follows time-respecting causal paths---a structure that static graph
representations erase. This thesis has examined whether graph neural network
architectures that explicitly encode temporal information can exploit this
structure to improve source detection accuracy, and under what conditions such
improvements are most pronounced.

\section{Summary of Contributions}
\label{st:summary}

The first contribution of this thesis is a faithful re-implementation of the
Backtracking Network (BN) of \cite{Ru2023} within a unified training and
evaluation framework based on Monte Carlo SIR simulation. The original BN was
demonstrated on a specific class of contact networks; our re-implementation
generalises the edge texture encoding to arbitrary temporal networks and
integrates it with the wandb-tracked experimental pipeline described in
Chapter~\ref{c:chapter5}, facilitating direct comparison against other methods
under identical data conditions.

The second contribution is the adaptation of the De Bruijn Graph Neural Network
(DBGNN) \cite{Qarkaxhija2022} and the DAG-GNN on Temporal Event Graphs for the
source detection task. Both architectures were originally proposed for
classification or prediction problems on temporal networks and required
non-trivial modification---in particular, the replacement of their output heads
and loss functions with a node-ranking formulation appropriate for source
detection, and the development of dedicated data preprocessing pipelines
(De Bruijn graph construction and temporal event graph construction,
respectively) that interface with the shared \texttt{TSIRData} dataclass.

The third and central contribution is the first systematic multi-architecture
comparison of temporal GNN methods for epidemic source detection. By evaluating
the static GNN, BN, DBGNN, and DAG-GNN under identical simulation budgets,
network families, and evaluation metrics---top-$k$ accuracy, average error
distance, and reciprocal rank---we provide a principled empirical basis for
architectural choice in this domain. The comparison also includes classical
heuristics (Jordan center, degree centrality, uniform random, and the
Soft-Margin Estimator baseline of \cite{Sterchi2025}), enabling assessment of
whether learned methods offer genuine advantages over analytical alternatives.

The fourth contribution is a robustness analysis structured around three
falsifiable hypotheses. H1 (Causal Path Integrity) predicts that temporal
architectures gain most in sparse, structured networks; H2 (Granularity
Convergence) predicts that their advantage diminishes as temporal resolution
decreases; and H3 (Temporal Signal-to-Noise Resilience) predicts that they
maintain accuracy at late observation times better than static models do. The
experimental results provide partial or full support for each hypothesis,
qualifying the conditions under which temporal GNNs constitute the preferred
approach and identifying the regimes in which the additional computational cost
of higher-order representations is justified.

\section{Limitations}
\label{st:limitations_ch}

The primary limitation of this work is that all results are obtained from
synthetic SIR simulations. Although the SIR model captures the essential
compartmental dynamics of many real infectious diseases, it abstracts away
heterogeneous infectivity, latency periods, reinfection, and the imperfect
contact-rate assumptions inherent in any recorded contact dataset. The models
trained and evaluated here may therefore exhibit different performance
characteristics when applied to real-world outbreaks where the true spreading
parameters are unknown and the observed infection states are subject to
reporting delays or under-ascertainment. Evaluating on real outbreaks with
known ground truth---such as well-documented laboratory or household
transmission studies---remains an important direction for future validation.

A further limitation concerns the scalability of the more expressive temporal
architectures. The De Bruijn graph at order $k$ has a node set of size
$O(|E|^k)$ where $|E|$ is the number of temporal edges, and the Temporal Event
Graph can contain $O(|E|^2)$ directed links in the worst case; neither
architecture was tested on networks with more than a few thousand nodes or
tens of thousands of temporal contacts. The single-source assumption shared by
all evaluated models is also a practical constraint: real outbreaks frequently
involve multiple simultaneous or near-simultaneous introduction events, a
scenario that none of the architectures is designed to handle. Finally, while
we follow the evaluation protocol of \cite{Sterchi2025} to guard against
inflated comparisons with the Soft-Margin Estimator baseline, the comparison
remains limited by the simulation budget, and conclusions about the relative
merit of learned versus classical methods should be interpreted with appropriate
caution.

\section{Future Work}
\label{st:futurework_ch}

Several natural extensions follow from the limitations identified above. The
most immediate is the generalisation to multi-source detection, where the model
must identify a small set $\mathcal{S}^* \subset \mathcal{V}$ of patient-zero
nodes rather than a single one. This requires either a multi-label output
formulation or a sequential prediction scheme in which identified sources are
removed and the procedure is iterated---an approach analogous to peeling
algorithms in combinatorial optimisation. A related extension is inductive
generalisation: all models evaluated here are transductive in the sense that
they are trained and tested on the same network. Incorporating per-node memory
mechanisms along the lines of Temporal Graph Networks (TGN) \cite{Rossi2020}
would allow a model trained on one contact network to be applied to an unseen
network without retraining, a capability that is essential for deployment in
real-time outbreak response settings.

On the architectural side, the DBGNN can be strengthened by integrating
Hypergeometric Graph Ensembles (HYPA) \cite{vonPichowski2024}, which use
statistical hypothesis testing to identify higher-order temporal patterns that
are anomalous relative to a randomised null model. Retaining only statistically
significant De Bruijn edges as an inductive bias for the GNN would reduce noise
from spurious higher-order patterns and potentially improve generalisation
across network families. Beyond epidemic spreading, the framework developed
here is directly applicable to other information diffusion processes on temporal
networks, including rumour propagation in social media, malware transmission in
computer networks, and cascading failures in infrastructure systems. Adapting
the simulation engine and evaluation metrics to these domains would substantially
broaden the impact of the methodology. Finally, integrating active querying
strategies---whereby the model requests the infection status of additional
nodes to resolve ambiguity among top candidates \cite{Sterchi2025}---offers a
principled way to improve accuracy at reduced observation cost, and represents
a promising bridge between the passive inference paradigm studied here and
practical epidemiological investigation.
